% AUTHOR: Amlal El Mahrouss
% PURPOSE: An In-Between Proof Before Algebraic Concepts.

\documentclass[10pt]{article}

%% ---- Core Packages ----------------------------------------
\usepackage[T1]{fontenc}
\usepackage[utf8]{inputenc}
\usepackage{lmodern}
\usepackage{microtype}

%% ---- Page Geometry ----------------------------------------
\usepackage[
top=1in, bottom=1in,
left=0.75in, right=0.75in,
columnsep=0.25in
]{geometry}

%% ---- Mathematics ------------------------------------------
\usepackage{amsmath, amssymb, amsthm, amsfonts}
\usepackage{mathtools}
\usepackage{bm}
\usepackage{textcomp}
\usepackage{esint}

%% ---- Figures & Tables -------------------------------------
\usepackage{graphicx}
\usepackage{booktabs}
\usepackage{multirow}
\usepackage{array}
\usepackage{caption}
\usepackage{subcaption}
\usepackage{float}

%% ---- Code Listings ----------------------------------------
\usepackage{listings}
\usepackage{xcolor}
\usepackage{algorithm}
\usepackage{algpseudocode}

\lstdefinestyle{codestyle}{
	backgroundcolor=\color{gray!8},
	basicstyle=\ttfamily\footnotesize,
	breakatwhitespace=false,
	breaklines=true,
	captionpos=b,
	commentstyle=\color{green!50!black},
	keywordstyle=\color{blue}\bfseries,
	stringstyle=\color{orange!70!black},
	numberstyle=\tiny\color{gray},
	numbers=left,
	numbersep=5pt,
	showstringspaces=false,
	frame=single,
	rulecolor=\color{gray!40},
	tabsize=2
}
\lstset{style=codestyle}

%% ---- Hyperlinks -------------------------------------------
\usepackage[
colorlinks=true,
linkcolor=blue!70!black,
citecolor=green!50!black,
urlcolor=blue!60!black
]{hyperref}
\usepackage{url}

%% ---- Bibliography -----------------------------------------
\usepackage[numbers, sort&compress]{natbib}

%% ---- Miscellaneous ----------------------------------------
\usepackage{enumitem}
\usepackage{siunitx}
\usepackage{cleveref}

%% ---- Theorem Environments ---------------------------------
\theoremstyle{plain}
\newtheorem{theorem}{Theorem}[section]
\newtheorem{lemma}[theorem]{Lemma}
\newtheorem{corollary}[theorem]{Corollary}
\newtheorem{proposition}[theorem]{Proposition}

\theoremstyle{definition}
\newtheorem{definition}[theorem]{Definition}
\newtheorem{example}[theorem]{Example}

\theoremstyle{remark}
\newtheorem{remark}[theorem]{Remark}

%% ---- Custom Commands --------------------------------------
\newcommand{\R}{\mathbb{R}}
\newcommand{\N}{\mathbb{N}}
\newcommand{\Z}{\mathbb{Z}}
\newcommand{\E}{\mathbb{E}}
\newcommand{\Prob}{\mathbb{P}}
\newcommand{\norm}[1]{\left\lVert #1 \right\rVert}
\newcommand{\abs}[1]{\left\lvert #1 \right\rvert}
\newcommand{\inner}[2]{\left\langle #1,\, #2 \right\rangle}
\newcommand{\ie}{\textit{i.e.}\xspace}
\newcommand{\eg}{\textit{e.g.}\xspace}
\newcommand{\etal}{\textit{et al.}\xspace}

%% ---- Caption Formatting -----------------------------------
\captionsetup{
	font=small,
	labelfont=bf,
	format=hang,
	justification=justified
}

\renewcommand{\qedsymbol}{\ensuremath{\blacksquare}}

\usepackage{lettrine}

%% ---- Header / Footer --------------------------------------
\usepackage{fancyhdr}
\pagestyle{fancy}
\fancyhf{}
\fancyhead[R]{\small\thepage}
\renewcommand{\headrulewidth}{0.4pt}

%% ============================================================
%%  FRONT MATTER
%% ============================================================

\title{%
	\vspace{-1.5em}%
	\large\textbf{An In-Between Proof Before Algebraic Concepts.}%
}
\author{By Amlal El Mahrouss}
\date{\today}

%% ============================================================
\begin{document}
%% ============================================================

\maketitle
\tableofcontents
\newpage

% =============================================
\section{Synopsis:}
% =============================================

\lettrine[lines=1, lhang=0.1, loversize=0.2]{L}ET the following theorems, proof, and definitions:

% =============================================
\subsection{Definition Of $f(x)$ in $\mathbb{R}$:}
% =============================================

\begin{definition}
	Let us define the following function $f(a, qk)$ in $\mathbb{R}$:
	\begin{align*}
		f(a, qk) \in \mathbb{R}, \, f(a, qk) \to \infty.
	\end{align*}
	Where there is variables $\forall a, qk \in \mathbb{R}$.
\end{definition}

\subsection{Several Theorems:}

\begin{theorem}
\label{theo:first-zeta}
For all $x \in \mathbb{R}$ that for $\frac{\zeta(3)}{\pi}$ in $\mathbb{R}$.:
\begin{align*}
	x &\coloneq \frac{\zeta(3)}{\pi},\\
	\exists \pi^{1/12} &> x, \forall x \neq 0, \forall x \in \mathbb{R}.
\end{align*}
\end{theorem}
\begin{proof}
Per Theorem~\ref{theo:first-zeta}. Let the following in $\mathbb{R}$ by induction.:\\(i) By the property of $\pi^{1/12}$ being transcendental, by the previously stated inequality in $\mathbb{R}.$:
\begin{align*}
		\exists \pi^{1/12} > x, \forall x \neq 0, \forall x \in \mathbb{R}, \, x \coloneq \frac{\zeta(3)}{\pi}.
\end{align*}
(ii) This conpletes the proof.
\end{proof}
\begin{remark}
	\label{remark:sum-both}
	Combining both formulas below:
	\begin{align*}
		\pi^{1/12} - \sum_{n=1}^{p}\frac{\zeta(3)}{\pi} + \frac{\pi^{13/12}}{\zeta(3)}.
	\end{align*}
	One may explore the following equations, and identities in a domain $\mathbb{D}.$:
	\begin{align*}
		\pi^{1/12} - \sum_{n=1}^{p}\frac{\zeta(3)}{\pi} &= \mathcal{H}_n,\\
		\mathcal{P}_n &= \frac{\pi^{13/12}}{\zeta(3)}.
	\end{align*}
\end{remark}
\begin{remark}
\label{remark:frac-zeta}
	Building upon previous results (1.2, 1.3):
	\begin{equation}
		\mathcal{P} = \frac{\pi^{13/12}}{\zeta(3)} = \mathcal{P}_\phi
	\end{equation}
	We can also write the following generalization for $\mathcal{P}_k$ in $\mathbb{R}$:
	\begin{align*}
		\mathcal{P}_k &= \frac{\pi^{13/12}}{\zeta(k)}.
	\end{align*}
\end{remark}
\begin{remark}
	Building upon Remark~\ref{remark:frac-zeta}, there exists an inequality of:
	\begin{align*}
		\mathcal{P}_k > \mathcal{P}_\phi > \zeta(3) > \mathcal{P}_x. 
	\end{align*}
	Where $\mathcal{P}_k$ is a near approximation of $\zeta(3)$ in a Dirichlet series.
\end{remark} For $p \neq 0, p \in \mathbb{N}$. Noting the expression $\frac{\zeta(3)}{\pi}$, we will write the following theorems:
\begin{theorem}
There exist the following equation in $\mathbb{R}.$ And $\forall a, q, k \in \mathbb{R}.$ And $p \in \mathbb{N}.$:
\begin{align*}
	\sum_{k=1}^{p} \frac{f(a, qk)}{k} + (\pi^{1/12} - \frac{\zeta(3)}{\pi}) - \pi &= 0,
\end{align*}
Then moving the left-hand-side $\pi$ to a right-hand-side $\pm \pi$.:
\begin{align*}
	\sum_{k=1}^{p} \frac{f(a, qk)}{k} + (\pi^{1/12} - \frac{\zeta(3)}{\pi}) &= \pm \pi.
\end{align*}
And an alternate forms by dividing $\zeta(3)$ to the right-hand-side or by also moving $(\pi^{1/12} - \frac{1}{\pi})$ to the right-hand-side.:
\begin{align*}
	\sum_{k=1}^{p} \frac{f(a, qk)}{k} + (\pi^{1/12} - \frac{1}{\pi}) &= \pm \frac{\pi}{\zeta(3)},\\
	\sum_{k=1}^{p} \frac{f(a, qk)}{k} &= \pm \frac{\pi}{\zeta(3)} - (\pi^{1/12} - \frac{1}{\pi}) .
\end{align*}
\end{theorem}
\begin{remark}
For $\pm \pi \in \mathbb{R},$ the plus-minus sign was added as a caution measure against rigor.
\end{remark}
\begin{corollary}
Per Remark~\ref{remark:sum-both}, we define $\sum_{k=1}^{p} \frac{f(a, qk)}{k} = \pm \frac{\pi}{(\pi^{1/12} - \frac{\zeta(3)}{\pi})}$ to be less than $\pi$ in $\mathbb{R}$.
\end{corollary} Moving on to the next section, additional theorems will be stated.

% =============================================
\section{Additional theorems and more explorations:}
% =============================================

\subsection{Expanding Definitions:}
On expanding in more domains and several thoughts:

\begin{theorem}
\label{theo-1}
	For a set $\mathcal{S} \subset \mathbb{D}$ in a domain $\mathbb{D}$:\\
	A set $\mathcal{S}$ of $\mathbb{D}$ defined values, is defined to be used as a sequence of functions $F_n$ in $\mathbb{D}$.
	\begin{align*}
		\mathcal{S} &= \{F_n(a, qk), \dots\} \subset \mathbb{D}
	\end{align*}
	Where $\mathcal{S}$ is a set of functions $F_n$ with a symmetric opposite of $F_n$ in $\mathcal{S}$.
\end{theorem}
\begin{remark}
	And thus, per \ref{theo-1}, the domain $\mathbb{D}$ shall not be equivalent to $\mathbb{C}$. In which there may be applications in computational theory and formal methods.
\end{remark}
% =============================================
\section{Topology and More:}
% =============================================
A continuous topological group of parallel points may be defined as for domains $\mathbb{S} \land \mathbb{O}$:

\begin{equation}
	\mathcal{G}(\mathbb{S}_n, \mathbb{O}_n).
\end{equation}
Thus such groups may be used to describe parallel subgroups $\mathcal{B}_n$ of a group $\mathcal{G}$.

% =============================================
\section{Going Forward:}
% =============================================

\lettrine[lines=1, lhang=0.1, loversize=0.2]{O}UR next note will be covering algebraic properties and some of their identities.

\nocite{*}

\bibliographystyle{acm}
\bibliography{chambadal1981dictionnaire, pinter2010book, weil1967basic, A_Course_of_Pure_Mathematics}

%% ============================================================
\end{document}
%% ============================================================